\documentclass[exercice,a4]{thedocument}%
\usepackage{pstricks}%
\usepackage{tabularx}%
\usepackage{texmaths-colors}%
\usepackage{pst-plot}%
\startdatakeys
    \source{24GENMATMEAG1}
    \BO{2026}
    \cycle{4}
    \competencesDuSocle{CA,RA}
    \objectifsApprentissage{C4A1aa,C4A2ad,C4B2cb,C4D2aa,C4D3ac}
    \motsCles{fonction,graphique,translation,statistiques,trigonométrie}
\enddatakeys
\begin{document}%
    \def\up#1{\textsuperscript{#1}}%


\medskip

\emph{Cet exercice est un questionnaire à choix multiple (QCM).\\
Pour chaque question, trois réponses ({\rm A, B} ou {\rm C}) sont proposées.\\
Une seule réponse est exacte.\\
Recopier sur la copie le numéro de la question et la lettre correspondant à la réponse exacte.\\
Aucune justification n’est demandée.}

\begin{center}
\begin{tabularx}{\linewidth}{|m{8.5cm}|*{3}{>{\centering \arraybackslash}X|}}\hline
\textbf{Question}&\textbf{Réponse A} &\textbf{Réponse B} &\textbf{Réponse C}\\ \hline
\textbf{1.~}
On considère la fonction $f$ définie par $f(x) = 3x - 2$.

Quelle est l’image de $-4$ par cette fonction ?&$-14$&$-10$&$-3$\\ \hline
\textbf{2.~}
Combien vaut $(-5)^3$ ?&$-125$ &$-15$ &125\\ \hline
\textbf{3.~}
Quelle est l’image du point J par la translation qui transforme C en A ?

\begin{center}
\psset{unit=1cm}
\begin{pspicture}(3.1,2.2)
\psdots(0,0)(1,0)(2,0)(3,0)(0,1)(1,1)(2,1)(3,1)(0,2)(1,2)(2,2)(3,2)
\uput[ur](0,1){A} \uput[ur](0,2){B} \uput[ur](1,2){C} \uput[ur](1,1){D}
\uput[ur](1,0){E} \uput[ur](0,0){F} \uput[ur](2,2){G} \uput[ur](3,2){H}
\uput[ur](3,1){I} \uput[ur](2,1){J} \uput[ur](3,0){K} \uput[ur](2,0){L}
\end{pspicture}
\end{center}&H&E&D\\ \hline
\textbf{4.~}

\begin{minipage}{2.5cm}Quel est l’antécédent de 3 par la fonction $f$ ?
\end{minipage} \hfill
\begin{minipage}{4cm}
\psset{unit=0.5cm}
\begin{pspicture*}(-1,-3.5)(4,4)
\psgrid[gridlabels=0pt,gridwidth=0.15pt,subgriddiv=1]
\psaxes[linewidth=1.25pt,labelFontSize=\scriptstyle]{->}(0,0)(-1,-3)(4,4)
\psplot[plotpoints=500,linewidth=1.25pt,linecolor=blue]{-1}{3}{3 2 x mul sub}
\uput[ur](1,1){\blue $\mathcal{C}_f$}
\end{pspicture*}
\end{minipage}&3&$- 3$&0\\ \hline
\textbf{5.~}
On a mesuré les tailles, en m, de sept élèves :

\[1,46~;~1,65~;~1,6~;~1,72~;~1,7~;~1,67~;~1,75\]

Quelle est la médiane, en m, de ces tailles ?&1,72&1,67&1,65\\ \hline
\begin{minipage}{4cm}
\textbf{6.~} {\small Dans le triangle ABC rectangle en A ci-contre, qui n’est pas en vraie grandeur, quelle est la valeur de $\cos \alpha$ ?}
\end{minipage}\hfill
\begin{minipage}{4cm}
\psset{unit=0.8cm}
\begin{pspicture}(5.8,2.5)
\pspolygon(0.6,0.4)(5.1,0.4)(0.6,2.2)%ABC
\uput[dl](0.6,0.4){A} \uput[d](5.1,0.4){B} \uput[l](0.6,2.2){C}
\uput[l](0.6,1.3){3} \uput[d](2.85,0.4){4} \uput[ur](2.85,1.3){5} \uput[ul](4.6,0.28){\small $\alpha$}
\psframe(0.6,0.4)(0.85,0.65)\psarc(5.1,0.4){0.5}{160}{180}
\end{pspicture}
\end{minipage}&0,8&0,75&0,6\\ \hline
\end{tabularx}
%\end{enumerate}
\end{center}

\bigskip
\showthesource
\end{document}
